\documentclass[12pt]{article}
\usepackage[utf8]{inputenc}
\usepackage[T1]{fontenc}
\usepackage[margin=1in]{geometry}
\usepackage{setspace}
\usepackage{xcolor}
\usepackage[colorlinks=true, urlcolor=blue!60!black, linkcolor=blue!60!black]{hyperref}
\setlength{\parindent}{2em}
\setlength{\parskip}{1.25em}
\renewcommand{\baselinestretch}{1.0}
\begin{document}
\begin{center}
{\LARGE \textit{Troxel v. Granville}, 530 U.S. 57 (2000)\par}
\vspace{0.8em}
{\large Case Brief\par}
\end{center}
\vspace{1em}

\section{Memory Jogger}
In a fractured 4--2--3 decision, the Supreme Court held that a Washington trial court's application of a ``breathtakingly broad'' nonparental visitation statute to override a fit mother's limits on grandparent visitation violated her fundamental Fourteenth Amendment substantive due process right to direct the care, custody, and control of her children (\href{https://www.law.cornell.edu/supremecourt/text/99-138}{Cornell LII}; \href{https://www.courtlistener.com/opinion/1087668/troxel-v-granville/}{CourtListener}).

\section{Detailed Case Facts}
Tommie Granville and Brad Troxel shared a relationship that ended in June 1991. Although they never married, they had two daughters together, Isabelle and Natalie Troxel. After separating from Granville in 1991, Brad lived with his parents, Jenifer and Gary Troxel (the girls' paternal grandparents), and regularly brought Isabelle and Natalie to his parents' home for weekend visitation (\href{https://www.law.cornell.edu/supremecourt/text/99-138}{530 U.S. at 60}).

In May 1993, Brad committed suicide. For several months after their son's death, Jenifer and Gary Troxel continued to see Isabelle and Natalie on a regular basis. In October 1993, however, Granville informed the Troxels that she wished to limit their visitation to one short visit per month and special holidays (\href{https://www.law.cornell.edu/supremecourt/text/99-138}{530 U.S. at 60--61, 71}).

In December 1993, the Troxels filed a petition in the Washington Superior Court for Skagit County seeking court-ordered visitation with Isabelle and Natalie under two Washington statutes, Wash. Rev. Code \S\S\ 26.09.240 and 26.10.160(3) (1994), though only \S\ 26.10.160(3) remained at issue before the U.S. Supreme Court. Section 26.10.160(3) provided: ``Any person may petition the court for visitation rights at any time including, but not limited to, custody proceedings. The court may order visitation rights for any person when visitation may serve the best interest of the child whether or not there has been any change of circumstances'' (\href{https://www.law.cornell.edu/supremecourt/text/99-138}{530 U.S. at 61}).

At trial, the Troxels requested two weekends of overnight visitation per month and two weeks of visitation each summer. Granville never sought to cut off visitation entirely; instead, she asked the court to order one day of visitation per month with no overnight stay, plus participation in her family's holiday celebrations. Nobody alleged, and no court found, that Granville was an unfit parent (\href{https://www.law.cornell.edu/supremecourt/text/99-138}{530 U.S. at 61, 68, 71}).

In its oral ruling and 1995 decree, the Superior Court rejected Granville's proposal and ordered visitation for one weekend per month, one week during the summer, and four hours on each of the two petitioning grandparents' birthdays. The trial judge remarked that it was ``normally in the best interest of the children to spend quality time with the grandparent'' absent adverse lifestyle issues, effectively placing on Granville the burden to disprove that visitation was in her daughters' best interests, and drew on his own childhood memories of spending a week with his grandparents (\href{https://www.law.cornell.edu/supremecourt/text/99-138}{530 U.S. at 61, 69, 72}). While Granville's appeal was pending, she married Kelly Wynn, who had six other children; approximately nine months after the Superior Court entered written findings on remand, Wynn formally adopted Isabelle and Natalie (\href{https://www.law.cornell.edu/supremecourt/text/99-138}{530 U.S. at 61--62}).

\section{Procedural History}
\begin{itemize}
\item \textbf{Washington Superior Court for Skagit County (1995; remand findings):} Entered a visitation decree granting the Troxels one weekend per month, one week in the summer, and four hours on each grandparent's birthday. After an initial remand from the Court of Appeals for written findings, the Superior Court entered brief findings noting that the Troxels were part of a ``large, central, loving family'' able to offer ``opportunities for the children in the areas of cousins and music'' (\href{https://www.law.cornell.edu/supremecourt/text/99-138}{530 U.S. at 61--62}).
\item \textbf{Washington Court of Appeals (\textit{In re Troxel}, 87 Wash. App. 131, 940 P.2d 698 (1997)):} Reversed and dismissed the Troxels' petition on statutory grounds, holding that nonparents lacked standing under \S\ 26.10.160(3) unless a custody action was pending. Because it resolved the dispute on statutory standing, it did not expressly decide Granville's constitutional challenge (\href{https://www.law.cornell.edu/supremecourt/text/99-138}{530 U.S. at 62}).
\item \textbf{Washington Supreme Court (\textit{In re Smith}, 137 Wash. 2d 1, 969 P.2d 21 (1998)):} Consolidated the Troxels' case with two other nonparental visitation appeals and affirmed the dismissal in a 5--4 decision. Disapproving the Court of Appeals' statutory reading, it held that the plain text of \S\ 26.10.160(3) gave nonparents standing at any time, but struck the statute down on its face under the Federal Constitution for two reasons: (1) it required no threshold showing of harm or potential harm to the child, and (2) it swept too broadly by allowing ``any person'' to petition at ``any time'' under a bare best-interest standard (\href{https://www.law.cornell.edu/supremecourt/text/99-138}{530 U.S. at 62--63}).
\item \textbf{United States Supreme Court (530 U.S. 57 (2000)):} Granted certiorari (527 U.S. 1069 (1999)), heard oral argument on January 12, 2000, and on June 5, 2000 affirmed the judgment of the Washington Supreme Court by a 6--3 vote with no majority opinion (\href{https://tile.loc.gov/storage-services/service/ll/usrep/usrep530/usrep530057/usrep530057.pdf}{U.S. Reports PDF}).
\end{itemize}

\section{Judicial Votes}
The judgment affirming the Washington Supreme Court was 6--3, fractured 4--2--3 across six separate opinions (\href{https://www.law.cornell.edu/supremecourt/text/99-138}{530 U.S. at 59--60}):
\begin{itemize}
\item \textbf{Plurality (4):} Justice O'Connor announced the judgment of the Court and delivered an opinion joined by Chief Justice Rehnquist, Justice Ginsburg, and Justice Breyer.
\item \textbf{Concurring in the Judgment (2):} Justice Souter filed an opinion concurring in the judgment; Justice Thomas filed a separate opinion concurring in the judgment.
\item \textbf{Dissenting (3):} Justice Stevens, Justice Scalia, and Justice Kennedy each filed separate dissenting opinions.
\end{itemize}

\section{Holding}
\textbf{Judgment:} Affirmed. The judgment of the Washington Supreme Court invalidating the visitation order was affirmed (\href{https://www.law.cornell.edu/supremecourt/text/99-138}{530 U.S. at 75}).

\textbf{Rule Announced:} The Fourteenth Amendment's Due Process Clause protects the fundamental right of parents to make decisions concerning the care, custody, and control of their children. Because fit parents are presumed to act in their children's best interests, a court reviewing a third-party visitation petition over a fit parent's objection must accord at least ``some special weight'' to the parent's own determination. Applying a sweeping statute that lets ``any person'' petition ``at any time'' to override a fit parent's decision based solely on a judge's view of the child's best interests---while inverting the presumption against the parent and ignoring the parent's offer of meaningful visitation---violates substantive due process (\href{https://www.law.cornell.edu/supremecourt/text/99-138}{530 U.S. at 65--73}).

\textbf{Holding vs. Dicta:} Under \textit{Marks v. United States}, the narrowest ground supporting the judgment is Justice O'Connor's as-applied holding that the combination of Washington's unlimited statutory text, the trial court's failure to accord special weight to Granville's wishes (indeed, applying an opposite presumption favoring grandparents), Granville's uncontested fitness, and her willingness to permit some visitation rendered the order unconstitutional as applied. Crucial issues were expressly left undecided as non-holding or reserved questions: (1) whether the Due Process Clause requires all nonparental visitation statutes to mandate a threshold showing of harm or potential harm to the child; (2) the exact tier of scrutiny applicable to parental-rights infringements (the plurality avoided naming strict scrutiny, whereas Justice Thomas would have required it); and (3) whether more narrowly tailored state grandparent-visitation statutes are facially valid (\href{https://www.law.cornell.edu/supremecourt/text/99-138}{530 U.S. at 73}; \href{https://supreme.justia.com/cases/federal/us/530/57/}{Justia}).

\section{Analysis of Opinions}

\subsection{Plurality Opinion (Justice O'Connor, joined by Chief Justice Rehnquist and Justices Ginsburg and Breyer)}
Justice O'Connor began by acknowledging the demographic diversity of modern American families and the important caregiving role grandparents often play, while noting that third-party visitation statutes impose a substantial burden on the parent-child relationship. Tracing a 75-year lineage of substantive due process precedent---including \textit{Meyer v. Nebraska} (1923), \textit{Pierce v. Society of Sisters} (1925), \textit{Prince v. Massachusetts} (1944), \textit{Stanley v. Illinois} (1972), \textit{Wisconsin v. Yoder} (1972), \textit{Parham v. J. R.} (1979), \textit{Santosky v. Kramer} (1982), and \textit{Washington v. Glucksberg} (1997)---she declared that parents' interest in the care, custody, and control of their children is ``perhaps the oldest of the fundamental liberty interests recognized by this Court'' (\href{https://www.law.cornell.edu/supremecourt/text/99-138}{530 U.S. at 63--66}).

Turning to Wash. Rev. Code \S\ 26.10.160(3), Justice O'Connor called the statute ``breathtakingly broad'' because it allowed any third party to haul a parent into court at any time and placed the best-interest decision exclusively in a judge's hands without requiring any deference to the parent. Four combined factors made the order unconstitutional as applied: (1) Granville was indisputably a fit parent, triggering the \textit{Parham} presumption that fit parents act in their children's best interests; (2) the Superior Court gave no ``special weight'' to Granville's judgment and instead applied an inverted presumption favoring grandparent visitation; (3) the Superior Court gave no weight to the fact that Granville had already agreed to meaningful visitation; and (4) the court's slender findings (cousins, music, and the judge's personal childhood memories) showed a mere disagreement over what was ``better'' for the children. Declining to decide whether a showing of harm is constitutionally required or to announce a per se rule against all third-party visitation statutes, the plurality affirmed without remand to spare Granville further disruptive litigation (\href{https://www.law.cornell.edu/supremecourt/text/99-138}{530 U.S. at 67--75}).

\subsection{Concurrence in the Judgment (Justice Souter)}
Justice Souter would have affirmed on a narrower procedural footing: the Washington Supreme Court had authoritatively construed its own statute and invalidated it \textit{on its face} because it authorized any person at any time to seek court-ordered visitation subject solely to an unconstrained best-interest standard. Emphasizing that ``Meyer's repeatedly recognized right of upbringing would be a sham'' if a judge could override a parent's choice of a child's associates whenever the judge thought he could make a ``better'' decision, Justice Souter concluded that this facial overbreadth alone ended the case, avoiding any need to review the trial record as applied or to ``turn[] any fresh furrows in the `treacherous field' of substantive due process'' (\href{https://www.law.cornell.edu/supremecourt/text/99-138}{530 U.S. at 75--79}).

\subsection{Concurrence in the Judgment (Justice Thomas)}
Justice Thomas noted at the outset that neither party had challenged the Court's substantive due process precedents or argued that the original meaning of the Due Process Clause precludes judicial enforcement of unenumerated rights (and noted in a footnote that no Privileges or Immunities Clause claim was raised). Taking existing precedent such as \textit{Pierce v. Society of Sisters} as given, he criticized the plurality, Justice Souter, and Justice Kennedy for recognizing a fundamental parental right while curiously omitting the standard of review. Applying strict scrutiny, Justice Thomas concluded that Washington lacked ``even a legitimate governmental interest---to say nothing of a compelling one---in second-guessing a fit parent's decision regarding visitation with third parties'' (\href{https://www.law.cornell.edu/supremecourt/text/99-138}{530 U.S. at 80}).

\subsection{Dissenting Opinion (Justice Stevens)}
Justice Stevens argued first that certiorari should have been denied, and on the merits that the Court should reverse the Washington Supreme Court's facial invalidation and remand. Because \S\ 26.10.160(3) had a ``plainly legitimate sweep''---covering former custodial caregivers, de facto parents, and intimate family members---it could not be facially unconstitutional. Substantively, he rejected the Washington Supreme Court's rigid harm-to-the-child rule, emphasizing that visitation disputes are not a bipolar contest between parent and State: there is a third person whose interests are directly implicated, the child. Warning against treating children as ``so much chattel,'' Justice Stevens argued that children likely hold complementary liberty interests in preserving established family-like bonds, and that States must retain room to protect children from arbitrary parental decisions (\href{https://www.law.cornell.edu/supremecourt/text/99-138}{530 U.S. at 80--91}).

\subsection{Dissenting Opinion (Justice Scalia)}
Justice Scalia agreed as a matter of natural law that parents' right to direct their children's upbringing is among the ``unalienable Rights'' of the Declaration of Independence and the unenumerated rights ``retained by the people'' under the Ninth Amendment. Yet he insisted that neither document empowers federal judges to invalidate duly enacted state statutes. Treating \textit{Meyer}, \textit{Pierce}, and \textit{Yoder} as narrow relics of a largely repudiated substantive due process era, he refused to extend them to third-party visitation. Warning that judicial enforcement of unenumerated parental rights would inevitably usher in ``a new regime of judicially prescribed, and federally prescribed, family law,'' he argued that state legislatures---which are democratically accountable and able to correct mistakes quickly---should resolve such policy questions (\href{https://www.law.cornell.edu/supremecourt/text/99-138}{530 U.S. at 91--93}).

\subsection{Dissenting Opinion (Justice Kennedy)}
Justice Kennedy would have vacated and remanded because the Washington Supreme Court erred in holding that the Constitution categorically bars the best-interest standard absent a showing of harm. He stressed that the conventional nuclear family is not the reality for many children and that a fit parent's right vis-\`{a}-vis a complete stranger differs markedly from her right vis-\`{a}-vis a de facto parent or long-term caregiver. Because 49 States then used some form of best-interest review in third-party visitation statutes, he argued that constitutional limits should be ``elaborated with care'' on a case-by-case basis by state courts---while candidly acknowledging that the sheer financial and emotional burden of defending domestic-relations litigation can itself unconstitutionally disrupt the parent-child relationship (\href{https://www.law.cornell.edu/supremecourt/text/99-138}{530 U.S. at 93--102}).

\section{Examples: Future Applications}
\begin{enumerate}
\item \textbf{Same-Side (Unconstitutional Interference):} \textit{A State Enacts a ``Community Mentors'' Visitation Law.} A state statute allows any adult neighbor, former teacher, or family friend to petition for monthly overnight visits with a child whenever a judge finds the visits would enrich the child's social development, without requiring proof of parental unfitness or deference to the parents' objection. This falls squarely on the same side as \textit{Troxel}: like Washington's \S\ 26.10.160(3), it allows any third party to override fit parents' childrearing decisions based solely on a judge's independent view of what is ``better'' for the child.
\item \textbf{Same-Side (Unconstitutional Interference):} \textit{Trial Court Presumes Grandparent Visitation After a Parent Limits Visits.} Under a state grandparent-visitation statute that is silent on the burden of proof, a fit widowed father permits his deceased wife's parents to visit his son every other month, but the grandparents sue for weekly overnights. The trial judge orders weekly visits, stating that ``grandparents are presumptively good for kids unless the father proves otherwise.'' Even though the statute limits petitioners to grandparents, the trial court's inverted presumption and refusal to give special weight to a fit parent's reasonable offer of visitation violate \textit{Troxel} as applied.
\item \textbf{Opposite-Side (Constitutional State Action):} \textit{Clear-and-Convincing Harm Standard After Complete Cutoff.} A narrowly drafted state statute permits grandparents to petition for visitation only if a parent has died and the surviving parent has completely severed a longstanding, substantial relationship with the child; the statute codifies a rebuttable presumption that the fit parent's decision is in the child's best interest and allows court-ordered visitation only if the grandparents prove by clear and convincing evidence that severing the bond will cause significant psychological harm to the child. This falls on the opposite side of \textit{Troxel}: it restricts standing, accords strong ``special weight'' to the fit parent's decision, and satisfies even the rigorous harm standard the plurality declined to mandate.
\item \textbf{Opposite-Side (Constitutional State Action):} \textit{State Intervention Following a Finding of Parental Unfitness.} A juvenile court adjudicates a mother unfit due to documented substance abuse and chronic neglect, places her child in temporary kinship care with the child's aunt, and orders structured visitation with the child's siblings and maternal grandmother over the mother's objection. Because the \textit{Troxel} presumption applies only ``so long as a parent adequately cares for his or her children (i.e., is fit),'' a judicial finding of unfitness removes the case from \textit{Troxel}'s core protection.
\item \textbf{Fence-Sitter (Close / Unresolved Case):} \textit{De Facto Parent Petition by a Former Long-Term Domestic Partner.} An unmarried couple raises a child together from birth for seven years, with the biological mother's full consent and encouragement of a parent-like bond, before splitting up; the biological mother (who is concededly fit) then cuts off all contact with her former partner. Under a state ``de facto parent'' statute, the former partner petitions for visitation without alleging harm, relying on a best-interest standard combined with a statutory factor giving ``due weight'' to the legal parent's wishes. This sits squarely on the fence: the plurality reserved whether harm is required and never defined ``special weight,'' Justice Thomas's strict scrutiny would likely side with the biological mother, whereas Justices Stevens and Kennedy expressly highlighted de facto parents as third parties whose established caregiving relationships may constitutionally qualify a parent's unilateral veto.
\end{enumerate}

\section{Critique}
\textbf{Originalist and Textualist Critiques:} Originalist scholars and jurists view \textit{Troxel} as an instructive fault line within conservative constitutional theory. Justice Scalia's dissent exemplifies judicial-restraint originalism: even if parental authority is a pre-political natural right, the Due Process Clause was originally understood as a guarantee of fair legal process rather than a license for federal judges to construct substantive family law (\href{https://www.law.cornell.edu/supremecourt/text/99-138}{530 U.S. at 91--93}; \href{https://scholarship.law.georgetown.edu/context/facpub/article/3587/viewcontent/An_Originalist_Theory_of_Due_Process_of_Law.pdf}{Georgetown Law Scholarship}). By contrast, commentators analyzing ``selective originalism'' note that Justice Thomas---despite questioning substantive due process and flagging the Privileges or Immunities Clause---applied strict scrutiny under existing precedent because neither party asked to overrule \textit{Meyer} or \textit{Pierce} (\href{https://texaslawreview.org/selective-originalism-and-judicial-role-morality/}{Texas Law Review}; \href{https://fedsoc.org/case/troxel-v-granville}{Federalist Society}).

\textbf{Doctrinal and Progressive Critiques:} Family-law scholars have criticized \textit{Troxel} from both doctrinal and child-centered angles. Emily Buss famously described parental rights after \textit{Troxel} as ``adrift in the middle,'' arguing that the plurality's refusal to articulate a clear standard of review or define ``special weight'' left lower courts with little usable guidance and invited costly case-by-case litigation (\href{https://chicagounbound.uchicago.edu/supremecourtrev/vol2000/iss1/7/}{Buss, \textit{Sup. Ct. Rev.} (2000)}). Similarly, John DeWitt Gregory and David D. Meyer observed that the fractured 4--2--3 split produced deep confusion across state courts over how much deference parents are owed and whether a harm showing is required (\href{https://scholarlycommons.law.hofstra.edu/cgi/viewcontent.cgi?referer=&httpsredir=1&article=1144&context=faculty_scholarship}{Gregory, \textit{The Detritus of Troxel}}; \href{https://lawecommons.luc.edu/cgi/viewcontent.cgi?referer=&httpsredir=1&article=1283&context=facpubs}{Woodhouse \& Coupet, \textit{Rutgers L.J.} (2001)}). From a progressive and child-welfare perspective, scholars such as Barbara Bennett Woodhouse, Sacha Coupet, and Solangel Maldonado have echoed Justices Stevens and Kennedy in warning that an overly rigid parental-autonomy model treats children's relational interests as secondary and can harm children in nontraditional, blended, or kinship-care households where quasi-parents have served as primary caregivers (\href{https://lawecommons.luc.edu/cgi/viewcontent.cgi?referer=&httpsredir=1&article=1283&context=facpubs}{Woodhouse \& Coupet (2001)}; \href{https://www.semanticscholar.org/paper/When-Father-(or-Mother)-Doesn\%27t-Know-Best\%3A-and-V.-Maldonado/850aa01eb77f26db562bdf0358f9a455385dfa6e}{Maldonado, \textit{Iowa L. Rev.} (2003)}).

\textbf{Underlying Competing Values:} At its core, \textit{Troxel} reflects an enduring tension among four values: (1) \textit{parental autonomy and family privacy}, which shields pluralistic childrearing choices from majoritarian or judicial standardization; (2) \textit{children's independent relational welfare}, which may suffer when a fit parent abruptly severs a child's deep bond with a grandparent or caregiver; (3) \textit{federalism}, given the states' historic primacy over domestic relations; and (4) \textit{judicial competence}, since both substantive due process balancing and open-ended ``best interest'' determinations vest wide discretion in judges.

\section{Key Quotations}
\begin{quote}
``The liberty interest at issue in this case---the interest of parents in the care, custody, and control of their children---is perhaps the oldest of the fundamental liberty interests recognized by this Court.''
\end{quote}
\noindent --- Justice O'Connor, plurality opinion (\href{https://www.law.cornell.edu/supremecourt/text/99-138}{530 U.S. at 65}).

\begin{quote}
``Accordingly, so long as a parent adequately cares for his or her children (i.e., is fit), there will normally be no reason for the State to inject itself into the private realm of the family to further question the ability of that parent to make the best decisions concerning the rearing of that parent's children.''
\end{quote}
\noindent --- Justice O'Connor, plurality opinion (\href{https://www.law.cornell.edu/supremecourt/text/99-138}{530 U.S. at 68--69}).

\begin{quote}
``As we have explained, the Due Process Clause does not permit a State to infringe on the fundamental right of parents to make childrearing decisions simply because a state judge believes a `better' decision could be made.''
\end{quote}
\noindent --- Justice O'Connor, plurality opinion (\href{https://www.law.cornell.edu/supremecourt/text/99-138}{530 U.S. at 72--73}).

\begin{quote}
``I would apply strict scrutiny to infringements of fundamental rights. Here, the State of Washington lacks even a legitimate governmental interest---to say nothing of a compelling one---in second-guessing a fit parent's decision regarding visitation with third parties.''
\end{quote}
\noindent --- Justice Thomas, concurring in the judgment (\href{https://www.law.cornell.edu/supremecourt/text/99-138}{530 U.S. at 80}).

\begin{quote}
``If we embrace this unenumerated right, I think it obvious---whether we affirm or reverse the judgment here, or remand as Justice Stevens or Justice Kennedy would do---that we will be ushering in a new regime of judicially prescribed, and federally prescribed, family law.''
\end{quote}
\noindent --- Justice Scalia, dissenting (\href{https://www.law.cornell.edu/supremecourt/text/99-138}{530 U.S. at 93}).


\end{document}
